% ============================================================
%  Use Cases – AI-Assisted Process Guidance Tool
%  Crane Maintenance Knowledge Management System
%  Master Thesis, RPTU / Fraunhofer IESE, 2025
% ============================================================

\section{Use Cases\label{sec:use_cases}}

This chapter describes the use cases implemented in the \textit{AI-Assisted Process Guidance}
prototype. The system supports crane maintenance engineers in structured, evidence-based fault
diagnosis by combining a Retrieval-Augmented Generation (RAG) pipeline with a conversational
AI agent, session tracking, and collaborative reporting.

The use cases are grouped into four functional areas that correspond to the four screens of
the application: (1)~authentication and user management, (2)~fault intake and session
initialisation, (3)~AI-guided diagnostic interaction, and (4)~knowledge reuse through the
crane dashboard.

% ------------------------------------------------------------
\subsection{Actors\label{sec:actors}}
% ------------------------------------------------------------

\begin{description}
    \item[Maintenance Engineer]
        The primary actor. A qualified crane maintenance engineer who uses the system to
        diagnose faults, record measurements, and access past fault reports. An engineer
        can create and interact with their own sessions and read all reports generated by
        any engineer.

    \item[AI Guidance System]
        The internal AI agent powered by Claude (\texttt{claude-sonnet-4-6}) via the
        Anthropic API. It acts as an automated diagnostic co-pilot that retrieves
        technical knowledge, asks targeted questions, interprets measurements, and
        generates structured fault reports.

    \item[Knowledge Base]
        A static collection of nine structured technical manual documents indexed into
        a ChromaDB vector store. The knowledge base is queried automatically on every
        conversational turn by the RAG retrieval subsystem.
\end{description}

% ------------------------------------------------------------
\subsection{UC-01: Register Account\label{sec:uc01}}
% ------------------------------------------------------------

\begin{description}
    \item[Actor] Maintenance Engineer
    \item[Precondition] The engineer does not yet have an account.
    \item[Trigger] The engineer navigates to the application and selects \textit{Create Account}.
    \item[Main Flow]
        \begin{enumerate}[noitemsep]
            \item The engineer enters their full name, a unique username, and a password.
            \item The system validates that the username is not already taken.
            \item The system hashes the password using \texttt{bcrypt} and stores the user
                  record in the SQLite database.
            \item A JWT access token (8-hour expiry) is issued and the engineer is redirected
                  to the Issue Intake Form.
        \end{enumerate}
    \item[Alternative Flow]
        If the username is already taken, the system displays an error message and the
        engineer is prompted to choose a different username.
    \item[Postcondition] A verified engineer account exists and the engineer is authenticated.
\end{description}

% ------------------------------------------------------------
\subsection{UC-02: Login\label{sec:uc02}}
% ------------------------------------------------------------

\begin{description}
    \item[Actor] Maintenance Engineer
    \item[Precondition] The engineer has a registered account.
    \item[Trigger] The engineer navigates to the application login screen.
    \item[Main Flow]
        \begin{enumerate}[noitemsep]
            \item The engineer enters their username and password.
            \item The system verifies the password against the stored bcrypt hash.
            \item A JWT token is issued and stored in the browser session state.
            \item The engineer is redirected to the Issue Intake Form.
        \end{enumerate}
    \item[Alternative Flow]
        If credentials are incorrect, an error is displayed and no token is issued.
    \item[Postcondition] The engineer is authenticated and can access all protected features.
\end{description}

% ------------------------------------------------------------
\subsection{UC-03: Submit Issue Intake Form\label{sec:uc03}}
% ------------------------------------------------------------

\begin{description}
    \item[Actor] Maintenance Engineer
    \item[Precondition] The engineer is authenticated.
    \item[Trigger] The engineer selects \textit{New Issue} from the navigation bar.
    \item[Main Flow]
        \begin{enumerate}[noitemsep]
            \item The engineer selects the crane type from a dropdown
                  (e.g.,~\textit{Demag EKKE 5t}, \textit{Liebherr Tower Crane}).
            \item The engineer selects the affected component from a dynamic dropdown
                  populated based on the crane selection
                  (e.g.,~\textit{Hoist Motor}, \textit{Hoist Brake}, \textit{Wire Rope}).
            \item The engineer describes the observed fault in a free-text field.
            \item Optionally, the engineer provides environmental conditions, recent
                  maintenance changes, and any error codes or warning indicators.
            \item Upon submission, the system creates a new troubleshooting session record
                  in the database and constructs a \textit{context snapshot} from the
                  intake fields.
            \item The context snapshot is passed to the AI Guidance System, which
                  retrieves relevant technical knowledge from the ChromaDB vector store
                  and generates an opening diagnostic message.
            \item The engineer is redirected to the AI Guidance Interface.
        \end{enumerate}
    \item[Postcondition]
        A new session is created, the AI has acknowledged the fault context, and
        the diagnostic conversation has begun without repeating information already
        provided in the form.
\end{description}

% ------------------------------------------------------------
\subsection{UC-04: Conduct AI-Guided Diagnostic Conversation\label{sec:uc04}}
% ------------------------------------------------------------

\begin{description}
    \item[Actor] Maintenance Engineer, AI Guidance System
    \item[Precondition] A troubleshooting session has been created (UC-03).
    \item[Trigger] The engineer reads the AI opening message and begins responding.
    \item[Main Flow]
        \begin{enumerate}[noitemsep]
            \item The engineer types an observation or answer in the chat input field
                  and submits it.
            \item The system stores the engineer's message and assembles the full
                  conversation history.
            \item The RAG subsystem constructs a search query from the engineer's
                  message combined with the session context. It retrieves up to five
                  relevant chunks from the knowledge base: three filtered by component
                  and two from general retrieval.
            \item The AI Guidance System receives the conversation history, retrieved
                  evidence, all recorded measurements, and the current session state.
            \item The AI responds with a targeted follow-up question or diagnostic
                  instruction, referencing specific values from the retrieved manual chunks.
            \item The AI embeds a structured \texttt{session\_update} JSON block in its
                  response, which the backend parses to update \texttt{completed\_steps},
                  \texttt{likely\_causes}, and \texttt{current\_hypothesis} in the database.
            \item The retrieved evidence is displayed in the \textit{Evidence} panel.
                  The updated session state is reflected in the \textit{Session State} panel.
        \end{enumerate}
    \item[Alternative Flow]
        If the AI identifies a safety-critical condition (e.g.,~brake not holding rated
        load), it flags the issue immediately and recommends crane shutdown before
        continuing the diagnostic sequence.
    \item[Postcondition]
        The session record is updated with new completed steps, likely causes, and the
        current working hypothesis. The conversation is persisted in the database.
\end{description}

% ------------------------------------------------------------
\subsection{UC-05: Record Measurements\label{sec:uc05}}
% ------------------------------------------------------------

\begin{description}
    \item[Actor] Maintenance Engineer
    \item[Precondition] An active troubleshooting session exists.
    \item[Trigger] The engineer takes physical measurements during the diagnostic process
                   and wishes to record them in the system.
    \item[Main Flow]
        \begin{enumerate}[noitemsep]
            \item The engineer opens the \textit{Measurements} tab in the right panel
                  of the guidance interface.
            \item The engineer enters one or more values from the following fields:
                  voltage~(V), current~(A), temperature~(\degree C), load~(kg),
                  brake gap~(mm), insulation resistance~(M\(\Omega\)), and vibration~(mm/s\,RMS).
            \item An optional free-text notes field is available for qualitative observations.
            \item Upon submission, the measurement record is stored in the database
                  and linked to the current session.
            \item All recorded measurements are injected into the AI system prompt on
                  subsequent conversational turns, enabling the AI to interpret the
                  values against the reference specifications in the knowledge base.
        \end{enumerate}
    \item[Postcondition]
        Measurements are persisted and visible in the panel. The AI will reference
        and interpret them in its next response.
\end{description}

% ------------------------------------------------------------
\subsection{UC-06: Generate Fault Report\label{sec:uc06}}
% ------------------------------------------------------------

\begin{description}
    \item[Actor] Maintenance Engineer, AI Guidance System
    \item[Precondition]
        A troubleshooting session has at least one completed conversational exchange.
    \item[Trigger]
        The engineer clicks \textit{Generate Report} in the guidance interface or
        from the dashboard session row.
    \item[Main Flow]
        \begin{enumerate}[noitemsep]
            \item The system collects the full conversation transcript and all
                  recorded measurements for the session.
            \item The AI Guidance System receives the complete session data and is
                  prompted to synthesise a structured JSON fault report.
            \item The AI generates the following report fields:
                \begin{itemize}[noitemsep]
                    \item \textbf{Issue summary}: one-sentence description of the reported fault.
                    \item \textbf{Steps taken}: ordered list of diagnostic actions performed.
                    \item \textbf{Root cause}: identified fault origin or
                          \textit{``undetermined''} if inconclusive.
                    \item \textbf{Diagnosis}: narrative explanation of findings (2--4 sentences).
                    \item \textbf{Recommendations}: prioritised list of corrective actions.
                    \item \textbf{Severity}: one of \texttt{critical}, \texttt{high},
                          \texttt{medium}, or \texttt{low}, based on the safety and
                          operational impact of the fault.
                    \item \textbf{Follow-up required}: boolean flag indicating whether
                          the crane requires further inspection before returning to service.
                \end{itemize}
            \item The report is stored in the database and linked to the session.
            \item The session status is updated to \texttt{completed}.
            \item The engineer is redirected to the dashboard where the report is
                  immediately visible.
        \end{enumerate}
    \item[Severity Classification Criteria]
        The AI assigns severity according to the following definitions:
        \begin{description}[noitemsep]
            \item[\texttt{critical}] Immediate safety risk; crane must not be operated.
            \item[\texttt{high}] Significant fault; operation should be suspended pending repair.
            \item[\texttt{medium}] Degraded operation; repair should be scheduled within 48 hours.
            \item[\texttt{low}] Minor issue; address at next planned maintenance interval.
        \end{description}
    \item[Postcondition]
        A structured, AI-synthesised fault report is stored and accessible to all
        engineers in the system.
\end{description}

% ------------------------------------------------------------
\subsection{UC-07: View Own Session History\label{sec:uc07}}
% ------------------------------------------------------------

\begin{description}
    \item[Actor] Maintenance Engineer
    \item[Precondition] The engineer is authenticated and has at least one session.
    \item[Trigger] The engineer navigates to \textit{Dashboard} and selects
                   \textit{My Sessions}.
    \item[Main Flow]
        \begin{enumerate}[noitemsep]
            \item The dashboard retrieves all sessions belonging to the logged-in engineer.
            \item Sessions are displayed in reverse chronological order with crane type,
                  component, problem description, status, date, and severity badge.
            \item The engineer can filter sessions by crane type, component, status,
                  and severity.
            \item Expanding a session row reveals the full fault report (if generated),
                  a \textit{Resume Session} or \textit{View / Resume} button, and a
                  \textit{Generate Report} button if no report exists.
            \item Summary statistics (total sessions, completed, reports generated,
                  follow-up required) and two charts (issues by component, issues by
                  severity) are displayed above the session list.
        \end{enumerate}
    \item[Postcondition] The engineer has a complete view of their diagnostic history.
\end{description}

% ------------------------------------------------------------
\subsection{UC-08: Resume an Active Session\label{sec:uc08}}
% ------------------------------------------------------------

\begin{description}
    \item[Actor] Maintenance Engineer
    \item[Precondition] An active (non-completed) session exists for the engineer.
    \item[Trigger] The engineer clicks \textit{Resume Session} or \textit{View / Resume}
                   on a session row in the dashboard.
    \item[Main Flow]
        \begin{enumerate}[noitemsep]
            \item The system retrieves the full message history and measurements
                  for the selected session from the database.
            \item The engineer is navigated to the AI Guidance Interface with the
                  complete conversation history restored.
            \item The engineer continues the diagnostic conversation from where it
                  was left off.
        \end{enumerate}
    \item[Postcondition]
        The session is resumed with full context; no information from the previous
        interaction is lost.
\end{description}

% ------------------------------------------------------------
\subsection{UC-09: View All Engineers' Reports (Knowledge Reuse)\label{sec:uc09}}
% ------------------------------------------------------------

\begin{description}
    \item[Actor] Maintenance Engineer
    \item[Precondition] The engineer is authenticated. At least one other engineer has
                        generated a report.
    \item[Trigger] The engineer navigates to \textit{Dashboard} and selects
                   \textit{All Engineers}.
    \item[Main Flow]
        \begin{enumerate}[noitemsep]
            \item The dashboard retrieves sessions from all engineers in the system.
            \item Each session entry displays the engineer name, crane type, component,
                  problem description, status, and severity badge.
            \item Expanding a session row shows the full fault report generated for
                  that session, including diagnosis, root cause, steps taken,
                  recommendations, and recorded measurements.
            \item The engineer cannot resume or modify another engineer's session;
                  the view is read-only for fault reports.
            \item Aggregate statistics and charts reflect the full dataset across
                  all engineers.
        \end{enumerate}
    \item[Postcondition]
        The engineer gains access to documented fault patterns and resolutions from
        colleagues, enabling knowledge reuse and organisational learning.
\end{description}

% ------------------------------------------------------------
\subsection{UC-10: Retrieve Technical Knowledge via RAG\label{sec:uc10}}
% ------------------------------------------------------------

\begin{description}
    \item[Actor] AI Guidance System (automated, triggered on every chat turn)
    \item[Precondition] The ChromaDB vector store is initialised and populated with
                        the nine knowledge base documents.
    \item[Trigger] An engineer submits a message in the diagnostic chat (UC-04).
    \item[Main Flow]
        \begin{enumerate}[noitemsep]
            \item The RAG subsystem constructs a composite search query from the
                  component name, problem description, and the engineer's latest message.
            \item A component-filtered query retrieves the three most semantically
                  relevant chunks from the document corresponding to the selected component.
            \item A general query retrieves two additional chunks from the full corpus
                  to capture cross-component context.
            \item Duplicate chunks are removed; the final evidence set (up to five chunks)
                  is ranked by cosine similarity score.
            \item The retrieved chunks are injected into the AI system prompt as
                  \textit{Retrieved Technical Knowledge} with source and relevance metadata.
            \item The chunks are also displayed in the \textit{Evidence} panel of the
                  guidance interface for transparency.
        \end{enumerate}
    \item[Postcondition]
        The AI response is grounded in component-specific manual content, with every
        recommendation traceable to a knowledge base source.
\end{description}

% ------------------------------------------------------------
\subsection{Use Case Summary\label{sec:uc_summary}}
% ------------------------------------------------------------

Table~\ref{tab:uc_summary} provides a consolidated overview of all ten use cases,
their primary actors, and the corresponding application screens.

\begin{longtable}{p{1.8cm} p{4.2cm} p{3.5cm} p{4cm}}
    \caption{Use case summary\label{tab:uc_summary}} \\
    \hline
    \textbf{UC ID} & \textbf{Name} & \textbf{Actor} & \textbf{Screen} \\
    \hline
    \endfirsthead
    \hline
    \textbf{UC ID} & \textbf{Name} & \textbf{Actor} & \textbf{Screen} \\
    \hline
    \endhead
    \hline
    \endfoot
    UC-01 & Register Account          & Engineer          & Login / Signup \\
    UC-02 & Login                     & Engineer          & Login / Signup \\
    UC-03 & Submit Issue Intake Form  & Engineer          & Issue Intake Form \\
    UC-04 & AI-Guided Diagnostic Chat & Engineer, AI      & AI Guidance Interface \\
    UC-05 & Record Measurements       & Engineer          & AI Guidance Interface \\
    UC-06 & Generate Fault Report     & Engineer, AI      & Guidance / Dashboard \\
    UC-07 & View Own Session History  & Engineer          & Crane Dashboard \\
    UC-08 & Resume Active Session     & Engineer          & Crane Dashboard \\
    UC-09 & View All Engineers' Reports & Engineer        & Crane Dashboard \\
    UC-10 & Retrieve Technical Knowledge (RAG) & AI System & Internal (all screens) \\
    \hline
\end{longtable}
